\documentclass[12 pt, letterpaper]{hmcpset}
\usepackage{hw-preamble}

\name{Jayden Thadani}
\class{MATH 142 --- Differential Geometry}
\assignment{Homework assignment \#4}
\duedate{23rd February 2024}

\begin{document}

\begin{problem}[Problem 1.]
	Let $[\mathbf{u}, \mathbf{p}], [\mathbf{v}, \mathbf{p}], [\mathbf{w}, \mathbf{p}] \in \mathrm{T}\mathbb{R}^{3}$ be tangent vectors. Prove the \textit{vector triple product identity}:
	\[
		[\mathbf{u}, \mathbf{p}] \times \left ( [\mathbf{v}, \mathbf{p}] \times [\mathbf{w}, \mathbf{p}] \right ) = \left ( [\mathbf{u}, \mathbf{p}] \cdot [\mathbf{w}, \mathbf{p}] \right ) [\mathbf{v}, \mathbf{p}] - \left ( [\mathbf{u}, \mathbf{p}] \cdot [\mathbf{v}, \mathbf{p}] \right ) [\mathbf{w}, \mathbf{p}].
	\]
\end{problem}

\begin{solution}
	Let $\mathbf{u} = (u_{1}, u_{2}, u_{3})), \mathbf{v} = (v_{1}, v_{2}, v_{3}), \mathbf{w} = (w_{1}, w_{2}, w_{3})$. Thus, we can compute
	\[
		\begin{split}
			[\mathbf{u}, \mathbf{p}] \times ([\mathbf{v}, \mathbf{p}] \times [\mathbf{w}, \mathbf{p}]) & = [\mathbf{u} \times (\mathbf{v} \times \mathbf{w}), \mathbf{p}]. \\
														   & = [\mathbf{u} \times (v_{2} w_{3} - v_{3} w_{3}, v_{3} w_{1} - v_{1} w_{3}, v_{1} w_{2} - v_{2} w_{1}), \mathbf{p}] \\
														   & = [(u_{2} v_{1} w_{2} - u_{2} v_{2} w_{1} - u_{3} v_{3} w_{1} + u_{3} v_{1} w_{3}, \\
														   & u_{3} v_{2} w_{3} - u_{3} v_{3} w_{2} - u_{1} v_{1} w_{2} + u_{1} v_{2} w_{1}, \\
														   & u_{1} v_{3} w_{1} - u_{1} v_{1} w_{3} - u_{2} v_{2} w_{3} + u_{2} v_{3} w_{2}), \mathbf{p}]
		\end{split}
	\]
	Notice that we can add to the $i$th component of the result the term $u_{i} v_{i} w_{i} - u_{i} v_{i} w_{i} = 0$  This allows us to get
	\[
		\begin{split}
			[\mathbf{u}, \mathbf{p}] \times ([\mathbf{v}, \mathbf{p}] \times [\mathbf{w}, \mathbf{p}]) & = [(u_{1} w_{1} + u_{2} w_{2} + u_{3} w_{3}) v_{1} - (u_{1} v_{1} + u_{2} v_{2} + u_{3} v_{3}) w_{1}, \\
														   & (u_{1} w_{1} + u_{2} w_{2} + u_{3} w_{3}) v_{2} - (u_{1} v_{1} + u_{2} v_{2} + u_{3} v_{3}) w_{2}, \\
														   & (u_{1} w_{1} + u_{2} w_{2} + u_{3} w_{3}) v_{3} - (u_{1} v_{1} + u_{2} v_{2} + u_{3} v_{3}) w_{3}), \mathbf{p}] \\
														   & = [(\mathbf{u} \cdot \mathbf{w}) \mathbf{v} - (\mathbf{u} \cdot \mathbf{v}) \mathbf{w}, \mathbf{p}] \\
														   & = \left ( [\mathbf{u}, \mathbf{p}] \cdot [\mathbf{w}, \mathbf{p}] \right ) [\mathbf{v}, \mathbf{p}] - \left ( [\mathbf{u}, \mathbf{p}] \cdot [\mathbf{v}, \mathbf{p}] \right ) [\mathbf{w}, \mathbf{p}].
		\end{split}
	\]
\end{solution}

\pagebreak

\begin{problem}[Problem 2.]
	\begin{enumerate}
		\item The \textit{arc length function} of a curve $\alpha : I \to \mathbb{R}^{3}$ is the map $s\left ( t \right ) = \int_{t_{0}}^{t} \left \Vert \alpha'\left ( u \right ) \right \Vert \, du$. Show that the curve $\alpha\left ( t \right ) = \left ( \cosh{t}, \sinh{t}, t \right )$ starting at $t_{0} = 0$ has arc length function $s\left ( t \right ) = \sqrt 2 \sinh{t}$.
		\item Say $t\left ( s \right )$ is the inverse of the arc length function $s\left ( t \right )$; then $\beta\left ( s \right ) = \alpha\left ( t\left ( s \right ) \right )$ is the \textit{unit-speed reparameterisation} of $\alpha$. Find the unit-speed parameterisation of $\alpha\left ( t \right ) = \left ( \cosh{t}, \sinh{t}, t \right )$
	\end{enumerate}
\end{problem}

\begin{solution}
	\begin{enumerate}
		\item First, to calculate the arc length parameterisation, we need to find the speed of the curve, $\Vert \alpha' \Vert$. Thus, we compute
			\[
				\begin{split}
					\Vert \alpha'(u) \Vert & = \Vert (\sinh{u}, \cosh{u}, 1) \Vert \\
							       & = \sqrt{\sinh^{2}{u} + \cosh^{2}{u} + 1} \\
							       & = \sqrt{\cosh(2u) + 1} \\
							       & = \sqrt{2 \cosh^{2}{u}} \\
							       & = \sqrt 2 \cosh u.
				\end{split}
			\]
			Note that the last step holds because $\cosh{u} \ge 0$ for all real $u$. 

			From here it's easy to compute the arc length --- we just integrate the speed from $0$ to $t$.
			\[
				\begin{split}
					s(t) & = \int_{0}^{t} \Vert \alpha'(u) \Vert \, du \\
					     & = \int_{0}^{t} \sqrt 2 \cosh{u} \, du \\
					     & = \sqrt 2 \left ( \sinh{u} \right )\Big |_{0}^{t} \\
					     & = \sqrt 2 \sinh{u}.
				\end{split}
			\]
	\item
		\[
			\boxed{
				\beta(s) = \left ( \cosh \left ( \sinh^{-1} \left ( {\frac{s}{\sqrt 2}} \right ) \right ), \frac{s}{\sqrt 2}, \sinh^{-1} \left ( \frac{s}{\sqrt 2} \right ) \right ).
			}
		\]

		Since $t$ is the inverse function of the smooth function $s$ which is given by $u \mapsto \sqrt 2 \sinh{u}$, then by the inverse function theorem, on the open interval $(0, \infty)$ (this excludes $0$ where $\sinh'{0} = 0$), we have
		\[
			t : u \mapsto \sinh^{-1}{\frac{u}{\sqrt 2}}
		\]
		as a smooth function. 

		Thus, we have the unit speed parameterisation $\beta$ given by $\alpha \circ t$ --- that is,
		\[
			\beta(s) = \left ( \cosh \left ( \sinh^{-1} \left ( {\frac{s}{\sqrt 2}} \right ) \right ), \frac{s}{\sqrt 2}, \sinh^{-1} \left ( \frac{s}{\sqrt 2} \right ) \right ).
		\]
	\end{enumerate}
\end{solution}

\pagebreak

\begin{problem}[Problem 3]
	Consider the curve $\alpha(t) = (2t, t^{2}, \log t)$ on the interval $I = \{ t \in \mathbb{R}, t > 0 \}$ .
	\begin{enumerate}
		\item Show that this curve passes through the points $\mathbf{p} = (2, 1, 0)$ and $\mathbf{q} = (4, 4, \log 2)$.
		\item The \textit{arc length} between two points $\alpha(t_{0})$ and $\alpha(t_{1})$ is the integral $\int_{t_{0}}^{t_{1}} \Vert \alpha'(u) \Vert \, du$. Find its arc length between the points $\mathbf{p}$ and $\mathbf{q}$.
	\end{enumerate}
\end{problem}

\begin{solution}
	\begin{enumerate}
		\item To show that the curve passes through $\mathbf{p}$  and $\mathbf{q}$ , just notice that
			 \[
				\begin{split}
					\alpha(1) & = (2 \cdot 1, 1^{2}, \log 1) \\
						  & = (2, 1, 0) \\
						  & = \mathbf{p}.
				\end{split}	
			\]
			and
			\[
				\begin{split}
					\alpha(2) & = (2 \cdot 2, 2^{2}, \log 2) \\
						  & = (4, 4, \log 2) \\
						  & = \mathbf{q}.
				\end{split}
			\]
			Since $\mathbf{p}$ and $\mathbf{q}$ are in the image of curve, the curve passes through them.
		\item 
			\[
				\boxed{
					\text{arc length} = 3 + \log 2
				}
			\] 

			By the definition of the arc length between two points, we have that
			\[
				\int_{1}^{2} \Vert \alpha'(u) \Vert \, du
			\]
			is the arc length between $\mathbf{p}$ and $\mathbf{q}$. To find the arc length explicitly, we compute $\Vert \alpha' \Vert$.
			\[
				\begin{split}
					\Vert \alpha'(u) \Vert & = \Vert (2, 2u, 1/u) \Vert \\
							       & = \sqrt{4 + 4u^{2} + 1/u^{2}} \\
							       & = \sqrt{\frac{4u^{4} + 4u^{2} + 1}{u^{2}}} \\
							       & = \sqrt{\frac{(2u^{2} + 1)^{2}}{u^{2}}} \\
							       & = \frac{2u^{2} + 1}{u}.
				\end{split}
			\]
			Then we get
			\[
				\begin{split}
					\int_{1}^{2} \Vert \alpha'(u) \Vert \, du & = \int_{1}^{2} 2u + \frac{1}{u} \, du \\
										  & = (u^{2} + \log(u)) \Big |_{1}^{2} \\
										  & = (4 + \log 2) - (1 + \log 1) \\
										  & = 3 + \log 2.
				\end{split}
			\]
	\end{enumerate}
\end{solution}

\pagebreak

\begin{problem}[Problem 4.]
	Suppose that $\beta_{1}$  and  $\beta_{2}$  are unit-speed reparameterisations of the same curve $\alpha : I \to \mathbb{R}^{3}$ .
	\begin{enumerate}
		\item Show that there is a number $s_{0}$ such that $\beta_{2}(s) = \beta_{1}(s + s_{0})$ for all $s$.
		\item What is the geometric significance of $s_{0}$?
	\end{enumerate}
\end{problem}

\begin{solution}
	\begin{enumerate}
		\item First we will show that any unit speed reparameterisation of $\alpha$ is an arc length parameterisation. Suppose $\beta : J \to \mathbb{R}^{3}$ is a unit speed reparameterisation of $\alpha$. Then note that we have $\beta = \alpha \circ h$ such that the following diagram commutes.
			\[
				\begin{tikzcd}
					J \ar[rd, "\beta"] \ar[d, "h"] \\
					I \ar[r, "\alpha"] & \mathbb{R}^{3}
				\end{tikzcd}
			\]
			We also have $\Vert \beta' \Vert = 1$ always, and $h'(s) \neq 0$ for all $s \in J$.

			By the chain rule, at any point $s \in J$, we have $\Vert \alpha'(h(s)) \, h'(s) \Vert = 1$. Thus, we have at every point $s \in J$
			\[
				\vert h'(s) \vert = \frac{1}{\Vert \alpha'(h(s)) \Vert}.
			\]
			Since we have $h' \neq 0$, either $h' > 0$ always or $h' < 0$ always. Assume without loss of generality that $h' > 0$ always, and thus,
			\[
				h'(s) = \frac{1}{\Vert \alpha'(h(s)) \Vert}
			\]
			for all $s \in J$. Since $\alpha$ is regular, this derivative is defined everywhere on $J$, and never $0$. Since $h$ is also regular, by the mean value theorem, there cannot be any distinct points in $I$ which take the same value under $h$ --- otherwise there would be a point between them where $h' = 0$. Thus, $h$ is injective. Since we have that $h$ is surjective (by definition), we have that $h$ is invertible, and, by the inverse function theorem, its inverse $\tau$ has derivative defined everywhere by
			\[
				\begin{split}
					\tau'(t) & = \frac{1}{h'(\tau(t))} \\
						 & = \Vert \alpha'(h(\tau(t))) \Vert \\
						 & = \Vert \alpha'(t) \Vert.
				\end{split}
			\]
			We can rewrite this as
			\[
				\tau(t) = \tau_{0} + \int_{t_{0}}^{t} \Vert \alpha'(u) \Vert \, du
			\]
			with $\tau(t_{0}) = \tau_{0}$ where $\tau : I \to J$. 

			In words, $\tau$ is an arc length function starting at $t_{0}$ and $h$ is its inverse. Thus, $\beta$ is an arc length parameterisation. 

			Now we know that when we have $\beta_{1}$ and $\beta_{2}$ as unit speed reparameterisations, they are defined by arc length parameterisations $\tau_{1}$ and $\tau_{2}$ with inverses $h_{1}$ and $h_{2}$. Since $\tau_{1}$ and $\tau_{2}$ are antiderivatives of the same function --- $\Vert \alpha' \Vert$, we have that $\tau_{2} = \tau_{1} - s_{0}$ for some $s_{0} \in \mathbb{R}$. 

			From this, we can get $h_{2}(s) = h_{1}(s + s_{0})$. Notice that for any $s \in J$ we have $s = \tau_{2}(t) = \tau_{1}(t) - s_{0}$. We can rewrite this as $\tau_{2}(t) = s$ and $\tau_{1}(t) = s + s_{0}$. Thus, taking $h_{1}$ and $h_{2}$ of both sides in the first and second equations respectively, we get
			\[
				h_{1}(s + s_{0}) = t = h_{2}(s)
			\]
			for any $s \in J$.
		\item $s_{0}$ is the arc length difference between the starting points of $\beta_{1}$ and $\beta_{2}$. That is, if $\beta_{1}$ starts at $t_{1}$ and $\beta_{2}$ starts at $t_{2}$, then $\alpha(t_{1})$ is an arc length of $s_{0}$ ``behind'' $\alpha(t_{2})$ $s_{0}$ (in the direction of $\beta_{1}$). Of course if $s_{0} < 0$ then, $\alpha(t_{1})$ is an arc length of $\vert s_{0} \vert$ ``in front of" $\alpha(t_{2})$.
	\end{enumerate}
\end{solution}

\pagebreak

\begin{problem}[Problem 5.]
	Let $\{ \mathbf{T}, \mathbf{N}, \mathbf{B} \}$  be the \textit{Frenet frame} associated with a unit-speed curve $\beta : J \to \mathbb{R}^{3}$ . That is, denote $\mathbf{T} = \beta'$  is the \textit{unit tangent vector}, $\mathbf{N} = \mathbf{T}'/\Vert \mathbf{T}' \Vert$  is the \textit{principal normal vector}, and $\mathbf{B} = \mathbf{T} \times \mathbf{N}$ is the \textit{binormal vector}. Prove the following identities:
	\begin{enumerate}
		\item $\mathbf{T} = \mathbf{N} \times \mathbf{B} = - \mathbf{B} \times \mathbf{N}$
		\item $\mathbf{N} = \mathbf{B} \times \mathbf{T} = - \mathbf{T} \times \mathbf{B}$
		\item $\mathbf{B} = \mathbf{T} \times \mathbf{N} = - \mathbf{N} \times \mathbf{T}$
	\end{enumerate}
\end{problem}

\begin{solution}
	\begin{enumerate}
		\item To show that $\mathbf{T} = \mathbf{N} \times \mathbf{B}$, consider $\mathbf{N} \times \mathbf{B}$. By the triple product identity we have
			\[
				\begin{split}
					\mathbf{N} \times \mathbf{B} & = \frac{\mathbf{T}'}{\Vert \mathbf{T}' \Vert} \times (\mathbf{T} \times \mathbf{N}) \\
								     & = \left ( \frac{\mathbf{T}'}{\Vert \mathbf{T}' \Vert} \cdot \mathbf{N} \right ) \mathbf{T} - \left ( \frac{\mathbf{T}'}{\Vert \mathbf{T}' \Vert} \cdot \mathbf{T} \right ) \mathbf{T} \\
								     & = \Vert \mathbf{N} \Vert^{2} \mathbf{T} - \frac{(\mathbf{T}' \cdot \mathbf{T}) \mathbf{T}}{\Vert \mathbf{T}' \Vert}.
				\end{split}
			\]
		Note that since $\mathbf{N} = \mathbf{T}'/\Vert \mathbf{T}' \Vert$ we have $\Vert \mathbf{N} \Vert = 1$ always, and similarly, since $\beta$ has unit speed $\Vert \beta' \Vert$ , we have $\Vert \mathbf{T} \Vert = 1$  always. Thus, thinking of $\mathbf{T}$  as a curve, by homework assignment \#3, problem 5, we have that $\mathbf{T} \cdot \mathbf{T}' = 0$  always. Substituting this into our previous work, we get
		\[
			\begin{split}
				\mathbf{N} \times \mathbf{B} & = 1 \mathbf{T} - \frac{0 \mathbf{T}}{\Vert \mathbf{T}' \Vert} \\
							     & = \mathbf{T}.
			\end{split}
		\]
		Then, by the anti-commutativity of the cross product $\mathbf{N} \times \mathbf{B} = - \mathbf{B} \times \mathbf{N}$ and thus, we have $\mathbf{T} = - \mathbf{B} \times \mathbf{N}$.
		\item To show that $\mathbf{N} = - \mathbf{T} \times \mathbf{B}$ we use the triple product identity again ---
			\[
				\begin{split}
					- \mathbf{T} \times \mathbf{B} & = - \mathbf{T} \times (\mathbf{T} \times \mathbf{N}) \\
								     & = - \left ( \mathbf{T} \cdot \frac{\mathbf{T}'}{\Vert \mathbf{T}' \Vert} \right ) \mathbf{T} + \left ( \mathbf{T} \cdot \mathbf{T} \right ) \mathbf{N}
				\end{split}
			\]
			Again, we have that $\mathbf{T} \cdot \mathbf{T} = \Vert \mathbf{T} \Vert = 1$ and thus, $\mathbf{T} \cdot \mathbf{T}' = 0$, and thus
			\[
				\begin{split}
					- \mathbf{T} \times \mathbf{B} & = - 0 \mathbf{T} + 1 \mathbf{N} \\
								       & = \mathbf{N}
				\end{split}
			\]
			Thus, by anti-commutativity we have $\mathbf{B} \times \mathbf{T} = - \mathbf{T} \times \mathbf{B}$ and thus, $\mathbf{N} = \mathbf{B} \times \mathbf{T}$.
		\item By definition, we have $\mathbf{B} = \mathbf{T} \times \mathbf{N}$ and by the anti-commutativity of the cross product we have $\mathbf{T} \times \mathbf{N} = - \mathbf{N} \times \mathbf{T}$, and thus, $\mathbf{B} = \mathbf{N} \times \mathbf{T}$.
	\end{enumerate}
\end{solution}

\end{document}
